\documentclass[a4paper,12pt]{article}
\usepackage[brazil]{babel}
\usepackage[utf8]{inputenc}
\usepackage[T1]{fontenc}
\usepackage{geometry}
\usepackage{setspace}
\usepackage{titlesec}
\usepackage{hyperref}
\usepackage{graphicx}
\usepackage{caption}
\usepackage{subcaption}
\usepackage{fancyhdr}
\usepackage{xcolor}

\input{commands.tex}

\hypersetup{
    colorlinks=true,% make the links colored
    linkcolor=blue
}

\usepackage{setspace}
\addbibresource{bibliography.bib}

\newcommand{\printingbibliography}{%

    \pagestyle{myheadings}
    \markright{}
    \sloppy
    \printbibliography[heading=bibintoc, % add to table of contents
                   title=Refer\^encias % Chapter name
                  ]
    \fussy%
}
\PassOptionsToPackage{table}{xcolor}
\renewcommand{\baselinestretch}{1.5}

\pagestyle{fancy}
\fancyhf{}
\renewcommand{\headrulewidth}{0pt}
\fancyhead[R]{\thepage}

\geometry{a4paper,top=30mm,bottom=20mm,left=30mm,right=20mm}

\titleformat*{\section}{\bfseries\large}
\titleformat*{\subsection}{\bfseries\normalsize}

\title{\Large Cinemática}
\author{Andr\'e V. Silva \\ \href{https://andrevsilva.com/}{andrevsilva.com}}
\date{\today}

\begin{document}

\maketitle

\tableofcontents

\newpage

A Cinemática é a parte da Mecânica Clássica que permite calcularmos a distância percorrida, 
a velocidade e a aceleração dos corpos.

Na Cinemática estudamos:
\begin{itemize}
    \item Movimento Retilíneo Uniforme (MRU);
    \item Movimento Retilíneo Uniformemente Variado (MRUV);
    \item Movimento Circular Uniforme (MCU);
    \item Movimento Circular Uniformemente Variado (MCUV).
\end{itemize}

Existem dois tipos de Cinemática:
\begin{itemize}
    \item \textbf{Cinemática Escalar}: não descreve o movimento em termos de direção e sentido, 
    apenas o módulo das grandezas.
    \item \textbf{Cinemática Vetorial}: descreve o movimento em termos de módulo, direção e sentido.
\end{itemize}

Cinemática e Cinética são termos que possuem significados diferentes. A energia cinética é a 
forma de energia relacionada ao movimento dos corpos.

Enquanto a Cinemática estuda o movimento dos corpos independentemente da sua causa, a Dinâmica 
estuda as causas do movimento dos corpos.

\textbf{O que se estuda em Cinemática?}

A Cinemática é o ramo da Mecânica Clássica que estuda o movimento dos corpos em trajetórias retilíneas 
e trajetórias circulares sem levar em consideração o que os provocou. Ela introduz diversos conceitos 
estudados na Física, como distância, velocidade e aceleração.

\section{\Large Fórmulas da Cinemática}

\subsection{Movimento Retilíneo Uniforme (MRU)}

No MRU, a velocidade é constante e a aceleração é nula ($a = 0$).

\subsubsection{Velocidade média}
\[
v_m = \frac{\Delta s}{\Delta t}
\]
Onde:
\begin{itemize}
    \item $\Delta s = s - s_0$ é o deslocamento
    \item $\Delta t = t - t_0$ é o intervalo de tempo
\end{itemize}

\subsubsection{Função horária da posição}
\[
s = s_0 + vt
\]
Onde:
\begin{itemize}
    \item $s$ é a posição final
    \item $s_0$ é a posição inicial
    \item $v$ é a velocidade constante
\end{itemize}

\subsection{Movimento Retilíneo Uniformemente Variado (MRUV)}

No MRUV, a aceleração é constante.

\subsubsection{Aceleração média}
\[
a = \frac{\Delta v}{\Delta t}
\]

\subsubsection{Função horária da velocidade}
\[
v = v_0 + at
\]

\subsubsection{Função horária da posição}
\[
s = s_0 + v_0 t + \frac{1}{2}at^2
\]

\subsubsection{Equação de Torricelli}
\[
v^2 = v_0^2 + 2a\Delta s
\]

\subsection{Queda Livre}

Movimento vertical sob ação exclusiva da gravidade ($g \approx 9{,}8\,m/s^2$).

\subsubsection{Velocidade}
\[
v = v_0 + gt
\]

\subsubsection{Altura (posição)}
\[
h = h_0 + v_0 t + \frac{1}{2}gt^2
\]

\subsubsection{Equação de Torricelli}
\[
v^2 = v_0^2 + 2g\Delta h
\]

\textbf{Observação:} O sinal de $g$ depende do referencial adotado.

\subsection{Lançamento Vertical}

Movimento vertical com subida e descida.

\subsubsection{Velocidade}
\[
v = v_0 - gt
\]

\subsubsection{Posição}
\[
s = s_0 + v_0 t - \frac{1}{2}gt^2
\]

\subsubsection{Equação de Torricelli}
\[
v^2 = v_0^2 - 2g\Delta s
\]

\textbf{Ponto importante:} No topo da trajetória, $v = 0$.

\subsection{Lançamento Horizontal}

Movimento composto: MRU na horizontal e MRUV na vertical.

\subsubsection{Posição horizontal}
\[
x = v_0 t
\]

\subsubsection{Movimento vertical}
\[
y = \frac{1}{2}gt^2
\]

\subsubsection{Velocidades}
\[
v_x = v_0 \quad ; \quad v_y = gt
\]

\subsection{Lançamento Oblíquo}

Movimento bidimensional com decomposição vetorial.

\subsubsection{Componentes da velocidade inicial}
\[
v_{0x} = v_0 \cos\theta \quad ; \quad v_{0y} = v_0 \sin\theta
\]

\subsubsection{Movimento horizontal}
\[
x = v_0 \cos\theta \cdot t
\]

\subsubsection{Movimento vertical}
\[
y = v_0 \sin\theta \cdot t - \frac{1}{2}gt^2
\]

\subsubsection{Velocidade vertical}
\[
v_y = v_0 \sin\theta - gt
\]

\subsubsection{Tempo de subida}
\[
t_s = \frac{v_0 \sin\theta}{g}
\]

\subsubsection{Altura máxima}
\[
h_{max} = \frac{(v_0 \sin\theta)^2}{2g}
\]

\subsubsection{Tempo total de voo}
\[
t = \frac{2 v_0 \sin\theta}{g}
\]

\subsubsection{Alcance horizontal}
\[
A = \frac{v_0^2 \sin(2\theta)}{g}
\]

\subsubsection{Relação do alcance}
\[
A = v_{0x} \cdot t
\]

\subsection{Movimento Circular Uniforme (MCU)}

Velocidade angular constante.

\subsubsection{Deslocamento angular}
\[
\Delta \theta = \theta - \theta_0
\]

\subsubsection{Velocidade angular}
\[
\omega = \frac{\Delta \theta}{\Delta t}
\]

\subsubsection{Função horária da posição}
\[
\theta = \theta_0 + \omega t
\]

\subsubsection{Relação entre grandezas lineares e angulares}
\[
v = \omega r
\]

\subsubsection{Aceleração centrípeta}
\[
a_c = \frac{v^2}{r} = \omega^2 r
\]

\subsection{Movimento Circular Uniformemente Variado (MCUV)}

Aceleração angular constante.

\subsubsection{Aceleração angular}
\[
\alpha = \frac{\Delta \omega}{\Delta t}
\]

\subsubsection{Função horária da velocidade angular}
\[
\omega = \omega_0 + \alpha t
\]

\subsubsection{Função horária da posição angular}
\[
\theta = \theta_0 + \omega_0 t + \frac{1}{2}\alpha t^2
\]

\subsubsection{Equação de Torricelli angular}
\[
\omega^2 = \omega_0^2 + 2\alpha \Delta \theta
\]

\section{Aceleração Centrípeta}

A aceleração centrípeta é a aceleração responsável por manter um corpo em movimento circular, sempre apontando para o centro da trajetória.

\[
a_{cp} = \frac{v^2}{R} = \omega^2 R
\]

Onde:
\begin{itemize}
    \item $a_{cp}$: aceleração centrípeta [$m/s^2$]
    \item $v$: velocidade linear [$m/s$]
    \item $R$: raio da trajetória [$m$]
    \item $\omega$: velocidade angular [$rad/s$]
\end{itemize}

\section{Tipos de Cinemática}

A Cinemática é separada em duas áreas:

\begin{itemize}
    \item \textbf{Cinemática Escalar}: não considera a análise vetorial (módulo, direção e sentido) do movimento.
    \item \textbf{Cinemática Vetorial}: considera módulo, direção e sentido do movimento.
\end{itemize}

\section{Tipos de Movimento na Cinemática}

\subsection{Movimento Retilíneo Uniforme (MRU)}

Também chamado de movimento uniforme, ocorre quando o corpo se move com velocidade constante em trajetória retilínea.

\begin{itemize}
    \item Velocidade constante
    \item Aceleração nula ($a = 0$)
\end{itemize}

\subsection{Movimento Retilíneo Uniformemente Variado (MRUV)}

Também chamado de movimento uniformemente variado, ocorre quando há aceleração constante.

\begin{itemize}
    \item Aceleração constante e diferente de zero
    \item Velocidade variável
\end{itemize}

\subsection{Movimento Circular Uniforme (MCU)}

Movimento em trajetória circular com velocidade constante.

\begin{itemize}
    \item Velocidade escalar constante
    \item Velocidade angular constante
\end{itemize}

\subsection{Movimento Circular Uniformemente Variado (MCUV)}

Movimento circular com aceleração angular constante.

\begin{itemize}
    \item Aceleração angular constante
    \item Velocidade angular variável
\end{itemize}

\section{Diferenças entre Cinemática e Cinética}

\begin{itemize}
    \item \textbf{Cinemática}: estuda o movimento dos corpos sem considerar suas causas.
    \item \textbf{Cinética}: refere-se à energia cinética, que é a energia associada ao movimento dos corpos.
\end{itemize}

\section{Cinemática x Dinâmica}

\begin{itemize}
    \item \textbf{Cinemática}: estuda o movimento (trajetórias, velocidades e acelerações).
    \item \textbf{Dinâmica}: estuda as causas do movimento (forças, leis de Newton, energia, impulso).
\end{itemize}

\section{Importância da Cinemática}

A Cinemática é fundamental para o estudo da Física, pois introduz conceitos essenciais para compreender o movimento dos corpos.

Algumas aplicações práticas incluem:

\begin{itemize}
    \item Lançamento e queda de corpos;
    \item Lançamento de bolas em esportes;
    \item Movimento de veículos ao acelerar;
    \item Movimento em curvas.
\end{itemize}

\section{Exercícios Resolvidos de Cinemática}

\subsection*{Exercício 1}
Um carro percorre 100 m em 5 s. Determine a velocidade média.

\[
\textcolor{blue}{v_m = \frac{\Delta s}{\Delta t} = \frac{100}{5} = 20\ \text{m/s}}
\]

\subsection*{Exercício 2}
Um móvel em MRU possui velocidade de 20 m/s e posição inicial de 10 m. Determine a posição após 5 s.

\[
\textcolor{blue}{s = s_0 + vt = 10 + 20 \cdot 5 = 110\ \text{m}}
\]

\subsection*{Exercício 3}
Um objeto parte do repouso e atinge 30 m/s em 10 s. Calcule a aceleração média.

\[
\textcolor{blue}{a = \frac{\Delta v}{\Delta t} = \frac{30 - 0}{10} = 3\ \text{m/s}^2}
\]

\subsection*{Exercício 4}
Um corpo em MRUV possui $v_0 = 5$ m/s e $a = 2$ m/s$^2$. Determine a velocidade após 4 s.

\[
\textcolor{blue}{v = v_0 + at = 5 + 2 \cdot 4 = 13\ \text{m/s}}
\]

\subsection*{Exercício 5}
Um móvel parte do repouso com aceleração constante de 3 m/s$^2$. Determine a posição após 6 s.

\[
\textcolor{blue}{s = \frac{1}{2}at^2 = \frac{1}{2} \cdot 3 \cdot 6^2 = 54\ \text{m}}
\]

\subsection*{Exercício 6}
Um carro a 25 m/s freia com desaceleração de 5 m/s$^2$. Determine a distância até parar.

\[
\textcolor{blue}{v^2 = v_0^2 + 2a\Delta s}
\]
\[
\textcolor{blue}{0 = 25^2 + 2(-5)\Delta s}
\]
\[
\textcolor{blue}{\Delta s = \frac{625}{10} = 62{,}5\ \text{m}}
\]

\subsection*{Exercício 7}
Uma bola é lançada verticalmente para cima com velocidade inicial de 20 m/s. Determine:
(a) o tempo de subida; (b) a altura máxima.

(a)
\[
\textcolor{blue}{t = \frac{v_0}{g} = \frac{20}{10} = 2\ \text{s}}
\]

(b)
\[
\textcolor{blue}{h = \frac{v_0^2}{2g} = \frac{400}{20} = 20\ \text{m}}
\]

\subsection*{Exercício 8}
Um projétil é lançado com velocidade inicial de 30 m/s formando um ângulo de $30^\circ$. Determine o alcance.

\[
\textcolor{blue}{A = \frac{30^2 \sin(60^\circ)}{10} \approx 77{,}94\ \text{m}}
\]

\subsection*{Exercício 9}
Um corpo realiza MCU com velocidade de 10 m/s em raio de 5 m. Determine a aceleração centrípeta.

\[
\textcolor{blue}{a_c = \frac{10^2}{5} = 20\ \text{m/s}^2}
\]

\subsection*{Exercício 10}
Um móvel em MCUV possui $\omega_0 = 2$ rad/s e $\alpha = 1$ rad/s$^2$. Determine a velocidade angular após 4 s.

\[
\textcolor{blue}{\omega = 2 + 1 \cdot 4 = 6\ \text{rad/s}}
\]

%%%%%%%% Bibliography 
% Os comandos para incluir as referências bibliográficas
%\printingbibliography

\end{document}
